\documentclass[11pt,a4paper]{article}
\usepackage[T1]{fontenc}
\usepackage[utf8]{inputenc}
\usepackage{lmodern,microtype,amsmath,amssymb,amsthm,mathtools,stmaryrd}
\usepackage[margin=24mm,headheight=15pt]{geometry}
\usepackage{booktabs,longtable,tabularx,array,enumitem,listings,xcolor}
\usepackage{fancyhdr,needspace,etoolbox,xurl,hyperref,bookmark}
\definecolor{ink}{HTML}{243C35}
\definecolor{accent}{HTML}{566A47}
\definecolor{wash}{HTML}{F1F4F0}
\definecolor{muted}{HTML}{53605B}
\hypersetup{colorlinks=true,linkcolor=ink,citecolor=ink,urlcolor=accent,
 pdftitle={Sorrel: Contract-Typed Proof Plans for Human-Scale Mathematics},
 pdfauthor={OpenAI, prepared for Vladimir Reshetnikov},bookmarksnumbered=true}
\urlstyle{same}
\pagestyle{fancy}\fancyhf{}
\fancyhead[L]{\small\sffamily Sorrel: contract-typed proof plans}
\fancyhead[R]{\small\sffamily Round 4}
\fancyfoot[C]{\thepage}
\renewcommand{\headrulewidth}{0.3pt}
\setlength{\parindent}{0pt}\setlength{\parskip}{0.43em}
\setlength{\emergencystretch}{2.5em}
\setcounter{tocdepth}{2}
\setlist{itemsep=0.2em,topsep=0.3em,leftmargin=*}
\newcolumntype{Y}{>{\raggedright\arraybackslash}X}
\newcolumntype{P}[1]{>{\raggedright\arraybackslash}p{#1}}
\newcommand{\xlongequal}[1]{\mathrel{\overset{#1}{=}}}
\newcommand{\code}[1]{\texttt{\detokenize{#1}}}
\newcommand{\N}{\mathbb N}\newcommand{\Z}{\mathbb Z}
\newcommand{\Q}{\mathbb Q}\newcommand{\R}{\mathbb R}\newcommand{\C}{\mathbb C}
\newcommand{\Prop}{\mathsf{Prop}}\newcommand{\Type}{\mathsf{Type}}
\newcommand{\Beh}{\mathsf{Beh}}\newcommand{\Min}{\operatorname{Min}}
\newcommand{\Supp}{\mathcal L}\newcommand{\Front}{\mathcal F}
\newcommand{\Eval}{\operatorname{eval}}\newcommand{\CheckCert}{\operatorname{check}}
\numberwithin{equation}{section}
\newtheorem{theorem}{Theorem}[section]
\newtheorem{proposition}[theorem]{Proposition}
\newtheorem{lemma}[theorem]{Lemma}
\theoremstyle{definition}\newtheorem{definition}[theorem]{Definition}
\theoremstyle{remark}\newtheorem{remark}[theorem]{Remark}
\lstset{basicstyle=\small\ttfamily,columns=fullflexible,keepspaces=true,
 breaklines=true,showstringspaces=false,backgroundcolor=\color{wash},
 frame=single,rulecolor=\color{wash},framesep=5pt,
 xleftmargin=5pt,xrightmargin=5pt,aboveskip=0.6em,belowskip=0.6em}
\BeforeBeginEnvironment{lstlisting}{\par\noindent\begin{minipage}{\linewidth}}
\AfterEndEnvironment{lstlisting}{\end{minipage}\par}
\pretocmd{\section}{\Needspace{6\baselineskip}}{}{}
\pretocmd{\subsection}{\Needspace{4\baselineskip}}{}{}
\begin{document}
\hypersetup{pageanchor=false}
\begin{titlepage}
\vspace*{12mm}
{\sffamily\large\color{accent}MATHEMATICAL PROOF-LANGUAGE DESIGN\par}
\vspace{19mm}
{\sffamily\bfseries\color{ink}\fontsize{44}{48}\selectfont Sorrel\par}
\vspace{6mm}
{\sffamily\bfseries\color{ink}\fontsize{24}{29}\selectfont Contract-Typed Proof Plans\\for Human-Scale Mathematics\par}
\vspace{9mm}
{\sffamily\Large Properties, residual obligations, and explicit evidence interfaces over Lean\par}
\vspace{13mm}\rule{\textwidth}{0.7pt}\vspace{4mm}
\textbf{The proposal.} Let authors write the mathematical route. Elaborate it
into a reusable proof plan whose type states exactly what remains to be
established. Keep alternative sufficient assumptions, not just one successful
proof. Discharge routine work with Lean, Leant, and checked computation without
changing the authored statement.
\vspace{7mm}
\textbf{The increment over round three.} A precise finite support calculus,
residual-interface extraction, a scope-safe integration design, and an executed
miniature frontend. A property list is useful; a typed explanation of how it
licenses the next mathematical step is more useful.
\vfill
Prepared for Vladimir Reshetnikov\par
Fourth iteration of the Brainstorm campaign\par
6 September 2026\par
\vspace{5mm}
{\small\color{muted}Evidence status: source-based design and written proofs;
31 executed finite-model tests; 7,680 exhaustive seed-subset/atom comparisons.
No complete Sorrel-to-Lean frontend, Lean compilation of the supplied specimen,
or productivity experiment is claimed.}
\end{titlepage}
\hypersetup{pageanchor=true}\pagenumbering{arabic}

\begin{abstract}
Sorrel is a proposed mathematical authoring language implemented first as an
extension of Lean's elaboration discipline, not its trusted conversion rules.
Its central object is a \emph{contract-typed proof plan}: a mathematical
argument with a fixed target, explicit construction choices, and an ordered
interface of residual proof obligations. Completed plans produce ordinary Lean
terms. Incomplete plans are conditional constructions, never silently published
as proofs of their conclusions.

The two round-three syntheses strongly support same-object property views,
behavioral contracts, bounded inference, and certificate-producing services,
but also warn that another surface syntax without an authoring experiment
would add little. Sorrel therefore concentrates on the missing operational
connection between these features. A finite, grounded implication engine
computes antichains of sufficient assumption sets and retains checked
conditional derivations. This supports alternative proof routes, precise
repair frontiers, and support-restricted methods. The paper proves soundness,
termination, relative completeness, and a residual-frontier theorem for that
finite calculus, while separating it from the harder problem of dependent
Lean elaboration.

Worked reconstructions study a compact-support Cauchy-transform formula in
ProveIt, exact integral Frey coefficients, sharp Fabius regularity, and a finite
tensor-family construction in the Fermat repository. They expose locality,
almost-everywhere reasoning, exact division, and dependent structure choices
that a useful type discipline must preserve. The proposal also specifies
behavior-aware Leant integration, a shared reification boundary for certified
computation, and a measured route from library interfaces to possible native
refinement types. An executable Python model and a generated, explicitly
uncompiled Lean implication specimen accompany the article. Their scope is
narrow and reproducible: they test the support calculus, not a finished proof
assistant.
\end{abstract}
\clearpage
{\small\setlength{\parskip}{0pt}\tableofcontents}
\clearpage

\section{A fourth-round decision, not a tenth spelling}\label{sec:decision}

\subsection{The central idea}
A mathematician's instruction ``apply the fundamental theorem of calculus''
usually identifies an argument more clearly than a sequence of coercion,
continuity, and interval-restriction steps. It does not eliminate those steps.
The interesting language problem is to make them \emph{generated consequences
of a stated mathematical route}, with a useful explanation when they fail.

Sorrel makes that route a typed object. A plan consists of a fixed mathematical
request, a sequence of author-selected steps, and the premises needed to
complete those steps. A plan may be reused under another context by supplying
its premises. Its intermediate objects remain available to later obligations.
Its completed evidence is an ordinary Lean term, together with a separately
checkable record of any required method policy.

For example, a quotient-transport plan should not merely report that it cannot
find a proof. It should report:
\begin{quote}
To identify this integer quotient with this rational quotient, establish exact
divisibility. The selected parity route needs oddness of the exponent. An
alternative registered route accepts a direct divisibility proof. The target
coefficient and its two coefficient domains have not changed.
\end{quote}
This is neither automatic hypothesis strengthening nor a claim that oddness is
mathematically necessary. It is an interface describing sufficient ways to
complete a particular argument.

\subsection{What is inherited and what is proposed here}
The round-three reports \emph{From Properties to Proof Obligations} and
\emph{Nine Refinements} already recommend fixed-object refinements, explicit
structures, dependent behavior specifications, bounded proof-producing
inference, exact-target replay, and an obligation frontier
\cite{synA,synB}. Sorrel adopts these decisions. Renaming a property row,
repeating the type of a hypothetical checker, or replacing \code{have} by a
new keyword would not be a substantive fourth iteration.

Sorrel's proposed increment is the following combination:
\begin{enumerate}
\item A proof plan exports its \emph{residual telescope}, not just a failure
message. The same object can drive author assistance, Leant completion, and
construction of a reusable helper theorem.
\item A precisely delimited inference layer retains \emph{alternative minimal
supports} within its registered rule fragment. This makes an explanation
stable under some edits and reveals alternative repairs without asserting
unrestricted weakest preconditions.
\item Method restrictions are checked over typed plan structure, independently
of proof irrelevance and kernel truth. Support policies are applied before
support minimization, so an unacceptable shortcut cannot suppress a permitted
argument.
\item A small executable model exercises the inference and residual-interface
mechanisms now. The full frontend and its evaluation remain explicit next
artifacts, not achievements inferred from the model.
\end{enumerate}

These are integration proposals relative to this campaign, not claims of
historical invention. The support calculus is closely related to
assumption-based truth maintenance and provenance annotations
\cite{dekleer,provenance}. The intended contribution is fitting such reasoning
into a dependent, proof-producing mathematical authoring interface without
confusing conditional support with established truth.

\subsection{A realistic definition of concision}
There are at least four different quantities: authored mathematical decisions,
authored formal bookkeeping, generated proof size, and reader effort. A system
can reduce one while increasing the others. A three-line call to a hundred-line
new helper theorem may be excellent library engineering; it is not by itself
evidence for a new language.

Sorrel therefore seeks to reduce repeated \emph{author interventions} while
preserving recoverable mathematical meaning. The selected domain, measure,
branch, witness, quantifier order, and strength of conclusion should remain
visible. Proof arguments for consequences already established by the context
may be folded. The expanded view must explain them using the actual accepted
plan, not a plausible narrative generated after the fact.

\section{Reading the corpus critically}\label{sec:corpus}

\subsection{Sources and the limits of this review}
Both main round-three TeX reports were examined, including their reconciliation,
worked examples, evidence qualifications, and fourth-iteration questions.
The design also uses direct readings of four mathematical source files and
Leant's verification interface. Table~\ref{tab:sources} fixes the source
identities. This is not an audit or rebuild of either mathematical repository,
nor a claim to have independently reread all nine original proposals.
Attribution to an individual round-three idea below is mediated by the two
syntheses unless otherwise indicated.

\begin{table}[htbp]\small
\caption{Source snapshots used in this article. Complete commit and blob
identities are recorded in \code{sources.json}.}\label{tab:sources}
\begin{tabularx}{\textwidth}{@{}P{28mm}P{24mm}Y@{}}\toprule
Repository & Commit prefix & Directly inspected material\\\midrule
Brainstorm & \code{bf333390} & Both round-three main TeX syntheses.\\
ProveIt & \code{24ce8bd7} & \code{CauchyCDF.lean}; \code{Regularity.lean}.\\
Fermat repository & \code{aa2d8b34} & \code{Def_FLTPrelim_FreyPackage.lean}; finite tensor-family solution.\\
Leant & \code{3a40904b} & \code{src/Leant/Synth/Verification.hs}; public synthesis and suggestion descriptions.\\\bottomrule
\end{tabularx}
\end{table}

The revised syntheses additionally discuss Leant at commit \code{823259f7}.
That account is useful, but it is not silently substituted for the
\code{3a40904b} reference resolved by the connector in this review. Claims
about the later behavioral-decision pipeline are attributed to the reports;
the callback-receipt distinction is independently visible in the inspected
Haskell source \cite{leantVerify,synA,synB}.

The reports describe successful Lean checks of seven earlier semantic cores,
including 49 named theorem declarations; the second gives 48 axiom inventories.
Those are their execution results, not fresh runs here. In particular,
Moraine's checked length-interpreter theorem must not be inflated into a
checked implementation of its coefficient normalizer. Schist's checked affine
example is a useful narrow end-to-end proof chain, not a general CAS. The
reports' deliberately failing \code{mvcgen} probe is not a clean whole-file
compilation \cite{synA,synB}.

\subsection{The important corrections must survive round four}
The synthesis process corrected several attractive but false simplifications.
Positive certificate soundness does not require decision completeness.
An abstract counterexample requires realization and the correct direction of
specification transfer. Over \code{List Empty}, different affine length
coefficients may describe the same concrete behavior. Seven similar Lean
encodings are not literally alpha-equivalent. And finiteness of an index set
does not imply that its algebra factors are finite modules
\cite{synA,synB}.

These corrections influence Sorrel's interfaces. A service advertises separate
positive, negative, and completeness contracts. A plan records its actual
input domain. A dependent construction exports the structures on its chosen
witness. An explanation says ``this registered route requires'' rather than
``the theorem requires'' unless necessity has separately been proved.

\subsection{Where actual Lean verbosity occurs}
The ProveIt Cauchy--CDF proof derives nonzero denominators from exclusion of a
complexified interval, obtains continuity and integrability, proves a primitive
identity on subintervals, transports it almost everywhere under a support
hypothesis, and applies a layer-cake identity. The argument is mathematically
coherent; much of its length connects these interfaces
\cite{cauchy}.

The Fabius regularity file similarly connects a differential equation,
nonnegative range bounds, real norms, nonnegative-real constants, and a
Lipschitz theorem. It separately proves optimality using a derivative value.
Its short corollaries are controls that Sorrel should leave short
\cite{regularity}.

The Frey file already uses a bundle containing mathematical properties.
Lean is not missing the ability to represent such information. The problem is
making operation-specific consequences available without repeatedly
unpacking and adapting the bundle. Its two displayed curve definitions also
show why identical notation across \(\Z\) and \(\Q\) can hide different division
semantics \cite{frey}.

Finally, the tensor-family proof constructs a type, ring and algebra
structures, finite/free proofs, a family of equivalences, and a naturality
law. These objects have dependent types: they cannot be represented faithfully
as an unordered bag of adjectives \cite{tensor}.

\subsection{The strongest existing baseline}
Mathlib's \code{fun_prop} already decomposes function properties, uses local
hypotheses, supports registered transition rules, and exposes side-condition
dischargers. Its documentation explicitly warns about the search cost of
transition rules \cite{funprop}. Sorrel must use and compare against this
capability rather than present continuity inference as a new feature.

Version discipline also matters. The corpus reports use Lean 4.32.0; the
August 10, 2026 release notes for Lean 4.33.0 describe further automation and
incremental elaboration changes, including the newer \code{vcgen} name for
\code{mvcgen'}. This does not retroactively change the pinned source experiment.
A modern comparison should specify its toolchain and include applicable
verification-condition generation, theorem search, and arithmetic tactics
\cite{lean433}.

\section{The author-facing language}\label{sec:surface}

\subsection{Write the argument, not an unbounded invitation to search}
Sorrel's mathematical surface is deliberately small. It has fixed variables,
assumptions, definitions, assertions, witness constructions, case splits,
induction, chains of equalities or inequalities, and named proof methods.
Ordinary Lean expressions and proofs remain escape hatches. Unrestricted
natural-language interpretation is not the acceptance boundary.

The following is \emph{proposed notation}, not syntax accepted by the supplied
finite-model parser:
\begin{lstlisting}
theorem reciprocal_integral (mu : FiniteMeasure Real)
    (a b : Real) (z : Complex)
    assuming a <= b, ae_support mu [a,b], z outside complexify [a,b]:
  integral x under mu of 1/(z-x)
    = mass(mu)/(z-b)
      - integral t from a to b of closed_cdf(mu,t)/(z-t)^2
proof
  let r(t) := 1/(z-t)
  let k(t) := r(t)^2
  have r'(t) = k(t) on [a,b] by differentiate
  have integral t from x to b of k(t) = r(b)-r(x)
    for x in [a,b] by fundamental_theorem
  use closed_cdf_layer_cake with mu, k
  replace the inner integral almost everywhere using the primitive
  conclude by linearity and additive algebra
\end{lstlisting}

The surface retains the actual mathematical decisions: the primitive,
layer-cake method, almost-everywhere replacement, and final rearrangement.
It delegates denominator nonzeroness, subinterval restrictions, and standard
continuity/integrability consequences. A proof of the whole theorem obtained
by unrelated search would not count as following these required steps.

The notation \code{FiniteMeasure Real} is an author-facing package convention.
At a Lean boundary it can elaborate to an ordinary measure and its
\code{IsFiniteMeasure} instance. The measure is not replaced, and the proof
of finiteness is not a second measure. Likewise, \code{closed_cdf} in this
finite-measure theorem denotes unnormalized cumulative mass, not a probability
CDF with an implicit normalization.

\subsection{Four verbs with different logical effects}
\code{assume P} extends a local proof context. It cannot be used at top level
to discharge a pre-existing obligation without making the resulting theorem
conditional. \code{have P by m} requests a proof of the fixed proposition
\(P\). \code{let x := e} fixes an object. \code{choose x satisfying S} delegates
a construction under the already elaborated specification \(S\).

This distinction prevents a common ambiguity: learning that an existing triple
is primitive is not the same action as dividing the triple by its gcd.
The latter constructs a different triple and owes a relation between the old
and new witnesses. Sorrel gives those actions different source nodes.

\subsection{Proof plans and holes}
A plan can be named before all its conditions are filled:
\begin{lstlisting}
plan cast_exact_quotient (n d q : Int)
  requires balance : d*q = n
           denominator : (d : Rat) != 0
  establishes (q : Rat) = (n : Rat)/(d : Rat)
  by map_balance_then_cancel
\end{lstlisting}
Here the input \(q\) is fixed. Another plan may construct it from \(d\mid n\),
with the zero-denominator policy explicit. The two plans are different APIs.
At use sites, \code{requires} arguments are proof-implicit: the elaborator
looks for evidence or leaves a typed goal. It does not ask a global instance
search to interpret every proposition as a typeclass.

A user can inspect a plan's residual interface, fill one obligation explicitly,
ask Leant to assemble a remaining adapter, or export the conditional plan as
a helper theorem. Exporting the conditional theorem does not establish an
unconditional target. The source view must keep that distinction visible.

\subsection{Hints, requirements, and publication}
A name written \code{by? arithmetic} is a search hint. A name written
\code{by arithmetic using only H} selects a required method and support
policy. This is an illustrative distinction; a concrete parser may choose
different punctuation, but the semantic difference is mandatory.

Every displayed assertion receives an evidence node. If the final theorem
happens not to use an earlier \code{have}, the earlier assertion must still be
checked before the document is marked complete. Unfinished assertions remain
visibly unfinished. A prose comment may explain intuition, but it is not
rendered as an established mathematical claim merely because it looks like
one.

\section{A property- and behavior-aware type discipline}\label{sec:types}

\subsection{Three contexts, one logical authority}
For explanation, distinguish a data context \(\Delta\), a proof context
\(\Gamma\), and a service profile \(\mathcal R\). The first fixes objects,
universes, operations, measures, domains, and any authorized construction
parameters. The second contains ordinary Lean hypotheses and accepted proofs.
The third selects registered theorem schemas, automation budgets, and method
policies. In an implementation, \(\Delta\) and \(\Gamma\) are classifications
of entries in Lean's dependent context, not competing logical environments.
Their dependency order must be retained.

A property view is written
\[
 \Delta;\Gamma\vdash t:A\mid\{P_1(t),\ldots,P_k(t)\}.
\]
Its meaning is simply that \(t:A\) is fixed and the context contains evidence
of each displayed proposition. The predicates may mention several objects;
for example, \(\mathsf{InversePair}(f,g)\), a naturality square, or
\(f=_{\mu\text{-a.e.}}g\). A row is a presentation of propositions, not a
restriction to unary properties or a new form of mathematical truth.

There are four elementary operations. Introduction adds \(p:P(t)\).
Intersection combines evidence about the same \(t\). Consequence applies an
explicit implication \(P(t)\to Q(t)\). Equality transport uses an actual proof
\(t=u\), or ordinary definitional equality, to obtain the corresponding
property of \(u\). None selects a different underlying multiplication,
topology, measure, or branch of a function.

\subsection{Property-implicit binders are a real source-level extension}
Consider a method of source type
\[
 \begin{aligned}
 &\mathsf{quotientDerivative}:\Pi f,g,f',g',x,\\
 &\quad \{h_f:\mathsf{HasDerivAt}(f,f',x)\}\to
 \{h_g:\mathsf{HasDerivAt}(g,g',x)\}\to\\
 &\quad \{h_0:g(x)\ne0\}\to T(f,g,x).
 \end{aligned}
\]
Braces here designate property-implicit parameters rather than ordinary
implicit data inference. At application, Sorrel checks the data arguments,
instantiates the exact required propositions, obtains evidence through the
selected profile, and inserts the proof arguments. When a proof is unavailable,
the application remains an open plan with that parameter in its interface.

This changes the author's typing discipline: uses of an operation are checked
against mathematical properties and behavior specifications. Yet the result
is an ordinary dependent application accepted by Lean. Calling it a type-system
extension is appropriate at the source language level; calling it a new Lean
kernel rule would be inaccurate.

\subsection{Local views versus exported types}
At an API boundary, a refined value can be represented by
\[
 \{x:A\mid P(x)\},
\]
or a structure with a value and proof fields. Lean's existing \code{Subtype}
already has this interpretation and is represented like the carrier in
compiled code. Proof irrelevance also avoids distinguishing subtype values
merely because their proof fields differ \cite{leanSubtype}.

Locally, repeatedly repackaging \(t\) is unnecessary. A view attaches evidence
to the already elaborated term. This retains convenient interaction with
existing APIs and avoids inventing chains of wrappers just to add more known
properties. The benefit to measure is fewer elaboration interventions, not
an unsupported claim that Lean's proof fields incur runtime storage costs.

Property polymorphism can be expressed without a new kind of row variable.
For instance, an interface parameter \(\rho:A\to\Prop\) can represent an open
family of caller-supplied properties. An operation that preserves that family
must have a theorem stating so. Sorrel must not infer arbitrary preservation
from a syntactically polymorphic type. In particular, a type parameter does
not by itself supply a parametricity theorem in classical Lean; the round-three
analyses correctly retain this warning \cite{synA,synB,leanAxioms}.

\subsection{Behavior is a relation involving the input and output}
For a fixed total function \(f:A\to B\), define
\[
 \Beh(f;P,Q):=\forall x:A,\ P(x)\to Q(x,f(x)),
 \qquad Q:A\to B\to\Prop.
\]
This represents, for example, a sorting function whose output is both ordered
and a permutation of \emph{its actual input}. A postcondition only about output
length is weaker. A constant list can satisfy length and orderedness while
losing the input elements; a constant empty list can satisfy orderedness alone.

Consequence is proof-producing:
\[
 \frac{\Beh(f;P,Q)\quad P'\Rightarrow P\quad
       \forall x,y,\ P'(x)\to Q(x,y)\to Q'(x,y)}
      {\Beh(f;P',Q')}.
\]
Thus a client supplying stronger premises can use an existing contract and
accept a weaker conclusion. Conversely, an implementation substituted for a
client's expected operation must demand no more and guarantee no less.
These are the two views of the same implication directions.

\begin{proposition}[Composition retains the intermediate]\label{prop:compose}
Let \(f:A\to B\), \(g:B\to C\), and suppose
\(\Beh(f;P,Q)\), \(\Beh(g;U,V)\). If
\begin{align*}
 &P(x)\to Q(x,y)\to U(y),\\
 &P(x)\to Q(x,y)\to V(y,z)\to W(x,z),
\end{align*}
then \(\Beh(g\circ f;P,W)\).
\end{proposition}
\begin{proof}
Fix \(x\) and \(h:P(x)\). The first contract gives \(Q(x,f(x))\).
The first bridge gives \(U(f(x))\); the second contract then gives
\(V(f(x),g(f(x)))\). Applying the second bridge proves the conclusion.
\end{proof}

A dependent operation may require the intermediate value, its index, or a
proof obtained at that stage. Its interface must keep those binders in order.
Flattening \(U(f(x))\) into a string called ``the next guard'' loses the
information needed for a well-typed application.

\subsection{Whole-function and higher-order constraints}
Many mathematical behaviors are not one-input/one-output postconditions.
Continuity, monotonicity, linearity, naturality, inverse-pair laws, and prefix
locality quantify over multiple evaluations or entire families. Sorrel accepts
arbitrary Lean propositions about the operation:
\[
 \mathsf{RelSpec}(f),\qquad
 \forall x,y,\ R(x,y)\to S(f(x),f(y)),\qquad
 \forall D,E,u,\ \mathsf{Commutes}(f_D,f_E,u).
\]
The inference profile supports only specified theorem shapes, but the type
language does not reduce these properties to output tags.

Higher-order composition follows the same rule. To synthesize an adapter on
invariant subtypes, the provider must prove that the original map preserves
the predicate. To lift a family to a natural transformation, the provider must
prove the commuting square. A matching arrow type provides neither theorem.

\subsection{Partial domains, totalized operations, and execution}
There is a semantic difference between
\[
 f:A\to B\quad\hbox{with}\quad\Beh(f;P,Q)
 \quad\text{and}\quad
 g:\{x:A\mid P(x)\}\to B.
\]
The first is total and has a conditional theorem; the second requires a
refined input. Neither licenses an arbitrary extension of \(g\) to all of
\(A\). Sorrel preserves imported Lean operations, including totalized division
and derivatives. A domain-sensitive mathematical API is an explicit new
interface, not a hidden reinterpretation of those operations.

Effectful programs need an operational semantics or an existing Hoare-style
specification framework. A stateful routine's contract may quantify over
initial and final states, exceptions, and termination. The surface can expose
these conditions, but the pure \(\Beh\) connective is not a substitute for
that semantics. Similarly, a stated cost bound must refer to a defined cost
model and be proved; a search timeout is an operational outcome, not a
mathematical precondition or complexity theorem.

\subsection{A native refinement former: worth considering, not assuming}
A later kernel experiment could introduce a primitive refinement former
\(\mathsf{Ref}(A,P)\) with explicit introduction \((a,p)\), erasing projection
\(\mathsf{val}\), a property projection, and proof-directed subsumption:
\[
 \frac{a:A\quad p:P(a)}{(a,p):\mathsf{Ref}(A,P)},\qquad
 \frac{u:\mathsf{Ref}(A,P)\quad h:\forall x,P(x)\to Q(x)}
      {\mathsf{weaken}(h,u):\mathsf{Ref}(A,Q)}.
\]
These rules have a straightforward candidate translation into existing Lean
subtypes. They might improve elaboration or transport handling in some
workloads, but they do not add the ability to express the properties discussed
here. Nor does the translation alone prove the implementation of a new kernel
sound.

A serious native design must specify universes, dependent substitution,
conversion, proof irrelevance, recursors, and the interaction with elimination
from \(\Prop\). It must also preserve the distinction between existence in
\(\Prop\) and executable data. Obtaining a witness through classical choice
is a different guarantee from constructing an executable algorithm.

The proposed decision gate is concrete: profile conservative Sorrel, isolate a
transport or checking bottleneck, implement the smallest native alternative,
and compare proof-term growth and checking cost. Independent kernel-checker
work such as Lean4Lean is relevant background for such an experiment
\cite{lean4lean}. Refinement calculi with carefully restricted checking rules
also show that the design space is not exhausted by unrestricted semantic
subtyping \cite{lovas}.

Sorrel rejects \emph{open-ended semantic entailment in conversion} as its
default. Inserting a checked implication application is predictable at the
kernel boundary; asking conversion to discover arbitrary implications is a
much stronger design commitment. CAS normalization likewise belongs in a
proof-producing step unless a separate, precisely specified extension has
its own metatheory and measurements.

\section{Residual interfaces: the type of an unfinished argument}\label{sec:plans}

\subsection{A residual telescope is not a set of suggestions}
An ordered telescope has the form
\[
 \Theta=(h_1:G_1,\ h_2:G_2(h_1),\ldots,
               h_k:G_k(h_1,\ldots,h_{k-1})).
\]
A selected plan elaborates to
\[
 \Delta;\Gamma\vdash\pi\rightsquigarrow
       [\Theta\ ;\ e:T],\qquad
 \Delta,\Gamma,\Theta\vdash_{\rm Lean}e:T.
\]
For a theorem request, \(T\) is the translation of the frozen authored
proposition. The residual telescope normally contains proof parameters.
Authorized data construction is represented explicitly by a construction
node or a dependent result type; it is not disguised as a routine proof hole.

Abstraction yields a conditional term
\[
 \lambda h_1\ldots h_k.e : \Pi\Theta,\ T.
\]
Supplying a well-typed substitution for \(\Theta\) yields a proof of \(T\).
Without that substitution, the plan proves only the conditional statement.
This is the exact boundary behind the user-facing word ``open.''

\subsection{Composition is typed substitution}
Suppose the next plan requires \(u:U\), and a previous plan constructs
\(e:U\) under residual telescope \(\Theta\). Replace \(u\) by \(e\) in the
next plan and its later obligations, append only dependencies that are
well-formed after that substitution, and recheck the resulting telescope.
Duplicate proposition obligations can be shared when their elaborated types
match; merely identical text is insufficient.

Consider an exact quotient. A construction may introduce \(q:\Z\) and a proof
\(dq=n\). The later cast step depends on \emph{that} \(q\). An obligation
about an independently synthesized \(q'\) cannot discharge it just because
both objects are described as ``the quotient.'' Either prove the appropriate
equality or retain the correct witness from the beginning.

\begin{proposition}[Conditional elaboration invariant]\label{prop:preservation}
For the restricted plan rules of property introduction, theorem application,
explicit transport, dependent substitution, and case/induction application,
assume that their Lean translations and data elaborations are well-typed.
If a plan has residual interface \([\Theta;e:T]\), then its conditional
translation has type \(\Pi\Theta,T\). Any well-typed filling of \(\Theta\)
produces a term of the same translated target \(T\).
\end{proposition}
\begin{proof}
Property introduction inserts its supplied proof. A method application uses
its registered theorem with the elaborated data and proof arguments.
Transport applies the supplied equality eliminator or relation-specific
consumer. Composition uses substitution into a well-formed dependent
context. Cases and induction apply their corresponding eliminators, with
all branch obligations abstracted in the scopes where they are valid.
Induction over the plan derivation gives a well-typed \(e\) under \(\Theta\);
repeated lambda introduction and substitution give the two conclusions.
\end{proof}

This is a paper invariant for the stated rules, not verification of a parser,
unifier, incremental elaborator, or arbitrary tactic. Its value is identifying
what such components must preserve, rather than claiming that kernel-checking
whatever they emit proves source fidelity.

\subsection{The semantic header is fixed before proof search}
The header contains the requested proposition, its free data parameters,
selected structures, quantifier order, and result relation. For example,
``equal almost everywhere under \(\mu\)'' and ``equal at \(x\)'' have different
headers. So do a bound and a least bound, or an arbitrary root and the root in
a specified interval.

Sorrel classifies unresolved metavariables before launching a provider.
Meaning-bearing data must be resolved from the source or from an explicitly
authorized construction specification. Proof metavariables remain obligations.
A provider is not allowed to instantiate an unknown measure, domain, or
coefficient structure merely because doing so makes the proof easy.

In a first Lean implementation, the header should be stored as the elaborated
expression with its context and universe instantiations. After provider
execution, the implementation checks that no protected assignment has changed.
A pretty-printed digest is useful for display and lookup but does not replace
the expression-level comparison. Source meaning is not preserved merely by
hashing the wrong expression consistently.

\subsection{Residual interfaces as reusable mathematical lemmas}
When an author repeatedly uses a plan, Sorrel can propose extracting a helper
whose parameters are the retained data and residual premises. This is a
mechanical interface extraction from an accepted conditional argument, not a
claim to have discovered the weakest possible theorem.

The Frey example motivates this sharply. A helper about \(a_4\) should not
inherit oddness and the residue of \(a\) when its chosen proof only uses
\(2\mid b\) and \(p\ge4\). Conversely, a theorem-level extraction must also
account for the dependencies of any called helper. The system cannot decide
support by searching for local hypothesis names in a pretty-printed proof.

A conservative implementation extracts the helper in a deliberately restricted
context and asks Lean to check it there. This establishes a reusable interface
with the permitted assumptions. It does not establish mathematical
indispensability of those assumptions, nor semantic consistency of the
context. A benchmark that is meant to exercise arithmetic should separately
supply an inhabitant of its premises.

\subsection{Why a method record is not just a proof term}
Lean's proof irrelevance makes proofs of one proposition logically
interchangeable; it is not a record of the author's intended argument
\cite{leanAxioms}. A proof term can be inspected operationally, but occurrence
of a theorem constant is not enough to show that a prescribed method actually
justified a step. A provider could insert a decorative use and then conclude
by an unrelated contradiction.

Sorrel's strict method mode therefore retains an \emph{intensional plan
certificate}. Its nodes apply approved theorem schemas to specified data and
premise nodes. The plan checker verifies the requested local conclusions,
allowed rules, and dependencies; Lean checks the resulting proof. A policy
may authorize arbitrary Lean proofs in designated leaves, or may require
those leaves to be closed in a restricted module. Those two policies are not
equivalent and must not share a misleading badge.

This does not solve the philosophical problem of recognizing the same human
argument in all equivalent proofs. It defines a narrow, inspectable guarantee:
the submitted plan belongs to this accepted grammar and supplies this fixed
target using these allowed theorem interfaces. A kernel-correct proof that
fails the method policy is reported as such, not called mathematically false.

\section{Alternative supports with a precise finite calculus}\label{sec:supports}

\subsection{The deliberately limited automatic fragment}
Fix a well-formed context and a finite universe \(\mathcal Q\) of ground
propositions. ``Ground'' means that the objects, structures, and other data
arguments have been fixed as elaborated expressions. Local variables such as
an arbitrary real \(x\) are allowed; there is no enumeration of real numbers.

Let \(H=\{H_1,\ldots,H_m\}\subseteq\mathcal Q\) be designated assumption
slots. They are possible \emph{inputs to a conditional plan}, not new axioms
accepted as true. A subset \(K\subseteq H\) records slots currently supplied
with actual evidence. Fix a finite rule table \(\mathcal R\), whose rules are
\[
 r:\ P_1\land\cdots\land P_k\Longrightarrow Q,
 \qquad P_i,Q\in\mathcal Q.
\]
In the intended Lean integration, every rule instance has a checked theorem
term. In the accompanying Python model, rules are stipulated implications,
so its certificates are conditional on that rule table.

\begin{definition}[Supports and labels]
A support of \(Q\) is a set \(S\subseteq H\) for which a finite rule derivation
of \(Q\) exists with assumption leaves drawn from \(S\).
The label \(\Supp(Q)\) is the collection of inclusion-minimal supports.
A retained witness records a finite derivation using exactly its leaf set.
\end{definition}

Minimality is inclusion-minimality inside this fixed vocabulary and rule
profile. It is not minimum cardinality, minimum effort, weakest semantic
precondition, or a claim that another library theorem cannot prove more.

\subsection{The algorithm}
Initialize \(\Supp(H_i)\) with \(\{H_i\}\), and every other label with the
empty family. A zero-premise rule contributes the empty support. For a rule
\(P_1,\ldots,P_k\Rightarrow Q\), combine retained supports of its premises:
\begin{equation}\label{eq:closure}
 \Supp(Q)\leftarrow\Min\left(
 \Supp(Q)\cup
 \left\{S_1\cup\cdots\cup S_k:
             S_i\in\Supp(P_i)\right\}\right),
\end{equation}
where \(\Min\) discards strict supersets. A newly inserted support receives
the proof tree formed from its premise witnesses and the actual rule.
Repeat until no label changes.

For instance, suppose the registered routes include
\[
 E\land L\Rightarrow D_{16},\qquad
 E\land L\land O\land A\Rightarrow D_4,\qquad
 D_4^{\rm direct}\Rightarrow D_4.
\]
For the joint goal \(D_4\land D_{16}\), the minimal supports are
\[
 \{E,L,O,A\},\qquad \{E,L,D_4^{\rm direct}\}.
\]
If \(E,L,A\) already have evidence, the open frontier is either \(O\) or a
direct proof of \(D_4\). Neither is installed as an assumption of the final
theorem without the author's explicit action.

\subsection{Soundness and termination}
\begin{theorem}[Support soundness]\label{thm:sound}
Every support retained by \eqref{eq:closure} has a finite derivation from exactly its recorded
assumption leaves. If the rules are interpreted by Lean proofs, that
derivation translates to a conditional Lean proof of its conclusion.
\end{theorem}
\begin{proof}
An initial support is witnessed by its own assumption. The empty support of a
zero-premise rule is witnessed by that rule. Every subsequent witness applies
a rule to already existing finite witnesses. Its leaf set is the union of
theirs, which is precisely the inserted support. Removing dominated supports
does not alter any retained witness. The Lean interpretation follows by
theorem application and abstraction over the leaves.
\end{proof}

\begin{theorem}[Finite termination]\label{thm:termination}
For fixed finite \(\mathcal Q,H,\mathcal R\), saturation terminates. At most
\(|\mathcal Q|2^m\) distinct support insertions can occur.
\end{theorem}
\begin{proof}
For each proposition, regard its label as representing the upward-closed
family of all supersets of its elements. An insertion strictly enlarges that
family; removing strict supersets during minimization changes no represented
support. A removed support can never be reinserted, since some retained subset
continues to dominate it, possibly through an even smaller replacement.
There are only \(2^m\) candidate subsets per proposition. Each complete pass
that does not terminate makes at least one such insertion.
\end{proof}

The bound counts successful label changes, not rule attempts or Lean cost.
Repeated Cartesian products, theorem instantiation, normalization, checking,
and witness construction are additional work. Even the number of incomparable
supports can be exponential: a family of all \(\lfloor m/2\rfloor\)-element
subsets has \(\binom m{\lfloor m/2\rfloor}\) members. This is a reason for
keeping the exact mode small, not for concealing the worst case.

\subsection{Relative completeness and schedule independence}
\begin{theorem}[Exactness of the final labels]\label{thm:complete}
After saturation, a set \(S\subseteq H\) derives \(Q\) by the fixed rule table
if and only if it contains some \(U\in\Supp(Q)\). Consequently the labels
are exactly the minimal supports and are independent of fair scheduling.
\end{theorem}
\begin{proof}
The reverse implication follows from Theorem~\ref{thm:sound} and weakening.
For the forward implication, use induction on the height of a finite
\(S\)-supported derivation. An assumption leaf has a retained support contained
in its singleton; an unconditional rule has empty support. For a rule node,
each premise derivation has a retained support contained in \(S\) by induction.
At saturation, their union either belongs to the conclusion's label or is
dominated by a retained smaller support. Thus the conclusion has a support
contained in \(S\). The collection of inclusion-minimal deriving sets is
uniquely determined by the rule table, so any saturated fair schedule yields
that same collection, although witness terms may differ.
\end{proof}

Unseeded cycles such as \(P\Rightarrow Q\Rightarrow P\) create no proof.
Once a seed exists, any retained proof still has a finite tree. This is not
coinductive reasoning. Deleting a seed and rerunning the algorithm cannot
justify the former result by its own cyclic dependency.

\subsection{The residual-frontier theorem}
Define
\[
 \Front_K(Q)=\Min\{S\setminus K:S\in\Supp(Q)\}.
\]
An empty element means a registered route is already supported. An empty
\emph{family} means no route has been derived in the fixed fragment. These
must have different displays.

\begin{theorem}[Sufficient repairs, exactly within the fragment]\label{thm:frontier}
For any set \(U\subseteq H\setminus K\) of additional supplied assumptions,
\(Q\) is derivable from \(K\cup U\) exactly when some
\(F\in\Front_K(Q)\) satisfies \(F\subseteq U\).
\end{theorem}
\begin{proof}
By Theorem~\ref{thm:complete}, derivability is equivalent to the existence of
\(S\in\Supp(Q)\) with \(S\subseteq K\cup U\). This is equivalent to
\(S\setminus K\subseteq U\). Removing dominated residual sets preserves
that existential condition.
\end{proof}

The theorem formalizes a useful editor operation: after removing a supplied
hypothesis, show the remaining ways to complete the same selected request.
A frontier element is a sufficient set of proof obligations, not permission
to assert them. If their proofs are hard or false in the current mathematical
context, the theorem remains open. The engine has not decided otherwise.

\subsection{Budgets, policies, and dependent boundaries}
A bounded run may retain only a sound subset of the final label families.
Its successful conditional proofs remain valid, but its frontier is
incomplete. An empty partial frontier proves neither impossibility nor absence
of another registered route. Sorrel reports the support inventory's
completeness separately from the existence of a checked proof.

Rules must be filtered for a required-method policy \emph{before} computing
labels, or labels must be indexed by a sufficiently precise policy state.
Suppose an unrestricted contradiction shortcut supplies a smaller support
than an allowed arithmetic argument. Discarding the arithmetic support first
and filtering afterward can erase the very route the user requested.
The supplied model reruns saturation with the allowed rule table; it does not
attempt a global policy automaton.

Finally, this calculus is not a solver for arbitrary dependent telescopes.
Its atoms live in one fixed, well-formed context. A premise involving a new
existential witness, a branch-local variable, or an unresolved structure cannot
be moved into the finite layer by giving it a string name. It first needs a
typed plan node that introduces the witness or scope. Support labels are local
to that node. This boundary is essential to the relevance of the finite
correctness results.

\section{Grounding, rewriting, and incremental scope}\label{sec:grounding}

\subsection{A finite rule set must actually be generated finitely}
Theorem~\ref{thm:termination} assumes a finite ground rule table. It does not
justify a frontend that first generates infinitely many theorem instances.
Sorrel's first implementation therefore uses an explicit \emph{grounding
profile}. Its inputs are the selected operation's signature, the local
context, a finite collection of eligible data terms, and a finite registry.

The initial term bank contains the request's subterms, named local objects,
and explicitly supplied construction results. A profile may add terms through
a whitelist of constructors to a fixed depth \(d\). If there are initially
\(M\) terms, at most \(C\) constructors of arity at most \(r\), and no fresh
unbounded literal generation, a crude finite bound is obtained from
\[
 T_0=M,\qquad T_{j+1}=T_j+C\max(1,T_j)^r.
\]
A schema with \(k\) eligible data parameters then has at most \(T_d^k\)
candidate assignments before type filtering. Universe instantiations and
structure arguments also come from explicitly finite choices.

Each assignment is checked against the theorem's actual dependent type.
Metadata labels binder roles; it does not override them. Proof arguments
become premise atoms. Schemas that require an unchosen data witness, higher-order
synthesis outside the permitted fragment, or a new unresolved structure are
not ground rules. They become explicit plan obligations or construction nodes.

This naive bound is intentionally pessimistic. A practical implementation uses
backward reachability from the demanded conclusion, discrimination indices,
and lazy generation. Its claim remains relative to the ground instances
actually covered by the profile. ``No supported instance generated'' is a
precise outcome; ``no proof exists'' is not.

\subsection{Reuse property automation instead of copying it}
The grounding layer should not replace \code{fun_prop}'s decomposition of
function expressions. A registered provider can solve a property goal using
\code{fun_prop} and return its proof as a leaf, while Sorrel tracks the
operation-level dependencies that connect it to other kinds of obligations.
A more explanatory provider may expose a validated decomposition into
subgoals. Both modes are useful, but only the latter supports fine-grained
frontier explanations through the provider's internal steps.

The same distinction applies to arithmetic. A successful \code{linarith}
call can close a guard without teaching Sorrel every arithmetic inference.
If the intended experiment is about explaining a removed inequality, the
provider must expose enough certified dependency information or be rerun in
restricted contexts. The extra explanatory work is a measurable cost.

\subsection{Anchors are typed expressions, not identifiers}
Suppose \(p:P(t)\) is available and simplification changes a use site from
\(t\) to \(u\). There are three legitimate cases. Definitional equality permits
ordinary reuse. A supplied equality \(e:t=u\) permits transport. A weaker
relation requires an appropriate consumer theorem. None is implemented
correctly by matching normalized strings.

Sometimes only the proposition changes: a proof of \(P(t)\leftrightarrow Q(u)\)
permits a proof of \(Q(u)\) without any claim that \(t=u\). With dependent
indices, both the term and the property family may need transport. Sorrel
asks Lean to check the resulting proposition, retaining universes and chosen
structure dictionaries. An unsupported transport shape remains an obligation.

The original fact need not disappear when a goal is rewritten. The index may
keep it at its original anchor and insert a checked bridge at the consumer.
This often avoids expensive global reindexing and preserves the explanation's
connection to the author's original object.

\subsection{Edits rebuild contexts; proofs do not migrate by hash}
A source edit can replace local identifiers and change the context in which a
term was elaborated. The initial Sorrel implementation gives the index the
lifetime of a Lean elaboration snapshot. A rebuilt declaration reconstructs
its index. Reuse across snapshots requires either re-elaboration from source
or a checked mapping of local declarations and their dependent types.

A source hash or receipt locates an artifact. It is not a proof that two
contexts are interchangeable. A renamed binder may be harmless; a different
measure hidden behind the same printed symbol is not. The safe initial
strategy is rebuilding, followed by selective cache improvements justified
by measurements and exact rechecking.

The finite model implements a much smaller invariant: every atom includes a
lexical context identifier, and a proof for one context is rejected in another.
It does not implement Lean expression transport, dependent index migration,
or a real editor. That difference is tested and reported rather than hidden
under the word ``scope-safe.''

\section{Worked reconstruction: an atom-exact Cauchy formula}\label{sec:cauchy}

\subsection{Statement and source correspondence}
Let \(\mu\) be a finite positive measure on \(\R\), \(a\le b\), and assume
\(\mu\)-almost every \(x\) lies in \([a,b]\). Fix
\(z\in\C\setminus\{t:t\in[a,b]\}\), with the canonical embedding
\(\R\hookrightarrow\C\) understood. Put
\[
 M=\mu(\R),\qquad F_\mu(t)=\mu(( -\infty,t]),\qquad r(t)=\frac1{z-t}.
\]
Here finite real masses are used; this is not an extended-real multiplication.
The directly inspected ProveIt theorem establishes
\begin{equation}\label{eq:cauchy}
 \int_\R\frac{d\mu(x)}{z-x}
 =\frac{M}{z-b}
   -\int_a^b\frac{F_\mu(t)}{(z-t)^2}\,dt.
\end{equation}
Its Lean implementation uses real scalar multiplication on complex values,
\code{Measure.real}, and oriented interval integrals \cite{cauchy}.
Those representation choices are part of the source correspondence, even when
a mathematical renderer uses conventional notation.

\subsection{The complete mathematical argument}
Define \(k(t)=(z-t)^{-2}\). Domain exclusion gives \(z-t\ne0\) on \([a,b]\).
Thus \(r\) and \(k\) are continuous there, and
\begin{equation}\label{eq:primitive}
 r'(t)=k(t),\qquad
 \int_x^b k(t)\,dt=r(b)-r(x)\quad(x\in[a,b]).
\end{equation}
Continuity on the compact interval and finiteness of \(\mu\), together with
its almost-everywhere support, imply \(\mu\)-integrability of \(r\).

The closed-CDF layer-cake identity is
\begin{equation}\label{eq:layercake}
 \int_a^b F_\mu(t)\,k(t)\,dt
 =\int_\R\left(\int_x^b k(t)\,dt\right)d\mu(x).
\end{equation}
One way to see it is to write
\(F_\mu(t)=\int\mathbf1_{\{x\le t\}}\,d\mu(x)\) and exchange the integrals.
On the supported interval the absolute value is bounded by a constant times
\(M(b-a)\), so the interchange is justified. For supported \(x\), the inner
integral runs from \(x\) to \(b\). Endpoint choices affect a Lebesgue-null set
in the \(t\)-integral and do not remove atoms of \(\mu\).

Substitute \eqref{eq:primitive} into the right side of \eqref{eq:layercake} almost everywhere:
\[
 \int_a^b F_\mu(t)k(t)\,dt
 =M r(b)-\int_\R r(x)\,d\mu(x).
\]
Rearranging gives \eqref{eq:cauchy}. The degenerate interval \(a=b\) is also
covered: its interval integral is zero and the measure is supported at that
point. The source proof realizes this argument through continuity,
interval-integrability, support restriction, the fundamental theorem, and
integral-congruence lemmas \cite{cauchy}.

\subsection{Which obligations should be implicit?}
The plan's domain layer proves
\[
 z\notin\operatorname{complexify}[a,b]\Longrightarrow
 \forall t\in[a,b],\ z-t\ne0.
\]
The regularity layer consumes that proposition to derive continuity and the
derivative formula. The integration layer uses the fixed measure, finiteness,
and support to derive integrability and almost-everywhere replacement.
The author supplies the layer-cake method and primitive identity; the
elaborator supplies well-typed applications of their standard premises.

These are different interfaces, not one universal ``regular'' flag. In
particular, an \(\mu\)-integrability fact is not a \(\nu\)-integrability fact,
and a derivative formula valid on \([a,b]\) must be restricted correctly to
\([x,b]\). Sorrel's plan must retain the interval-membership evidence needed
for each subinterval application.

A compact view can expose the following dependency summary without listing
all coercion lemmas:
\begin{center}\small
\begin{tabularx}{\textwidth}{@{}P{35mm}Y@{}}\toprule
Requested operation & Sufficient evidence consumed\\\midrule
Differentiate \(r\) & Nonzero denominator at the current point.\\
Use the primitive on \([x,b]\) & \(x\in[a,b]\); derivative on that subinterval; interval integrability.\\
Replace under \(\mu\) & Almost-everywhere support and the pointwise primitive for supported points.\\
Integrate \(r(b)-r(x)\) & Finiteness of \(\mu\) and \(\mu\)-integrability of \(r\).\\
Normalize mass to one & A separate probability-measure hypothesis.\\\bottomrule
\end{tabularx}
\end{center}

\subsection{Positive and negative neighbors}
For a point mass \(\mu=\delta_c\), \(c\in[a,b]\), the left side of
\eqref{eq:cauchy} is \((z-c)^{-1}\), while
\(F_\mu(t)=\mathbf1_{\{c\le t\}}\). The right side is
\[
 (z-b)^{-1}-\int_c^b(z-t)^{-2}\,dt=(z-c)^{-1}.
\]
This is an important positive fixture: no atomlessness premise is needed.
Requiring one would weaken the source theorem unnecessarily.

Removing domain exclusion should expose the nonzero-denominator requirements
of this route. It must not produce a universal claim that the identity is
false for all such inputs. For example, the zero measure can make a formula
trivial despite an unsuitable primitive route. Similarly, replacing finite
mass by unit mass without a probability hypothesis is a genuine change of
statement; a Dirac measure of mass two detects it immediately.

Almost-everywhere equality must not be promoted to endpoint equality. Under
Lebesgue measure, changing a function at one point preserves its integral but
can change its value at that endpoint. A consumer registry can use
\(f=_{\mu\text{-a.e.}}g\) to replace an integrand, while refusing point
application unless additional evidence is supplied.

\subsection{An explicit consumer registration for almost-everywhere rows}
A registered integral consumer has schematic type
\[
 (\forall^{\mu\text{-a.e.}}x,\ f(x)=g(x))
 \longrightarrow \int f\,d\mu=\int g\,d\mu.
\]
Other methods, such as distributing a difference into two integrals, have
additional integrability premises. The registry retains these exact theorem
interfaces rather than assuming that all integral operations share one guard.
For a family of almost-everywhere facts, a countable intersection may be
available under the appropriate theorem; an arbitrary uncountable intersection
is not inferred.

This answers a round-four question concretely: almost-everywhere information
enters as a relation with the measure in its anchor, and consumers are ordinary
Lean theorems with checked binder roles. It is not an ambient change from
pointwise mathematics to a quotient where all consumers become valid.

\section{Arithmetic and dependent construction controls}\label{sec:controls}

\subsection{Exact Frey quotients without a contradiction shortcut}
The Frey package contains nonzero integers \(a,b,c\), a prime exponent
\(p\ge5\), the Fermat equation, a gcd condition, and congruences. Its integral
and rational curve definitions both display
\[
 a_2=\frac{b^p-1-a^p}{4},\qquad
 a_4=\frac{-a^p b^p}{16},
\]
but division has different semantics in the two coefficient domains
\cite{frey}. The source bundle illustrates Lean's existing property-carrying
capability; Sorrel's task is extracting the premises that each use needs.

Following the corrected round-three benchmark, use instead the satisfiable
arithmetic context
\begin{equation}\label{eq:f0}
 a,b\in\Z,\quad p\in\N,\quad p\ge4,\quad p\text{ odd},\quad
 a\equiv3\pmod4,\quad 2\mid b.
\end{equation}
Since \(b=2u\), \(16\mid b^p\), and hence \(16\mid-a^pb^p\).
Also \(b^p\equiv0\pmod4\) and \(a^p\equiv(-1)^p=-1\pmod4\), so
\(4\mid b^p-1-a^p\). Only evenness and the exponent bound are needed for
\(a_4\); the other assumptions belong to the chosen \(a_2\) route.
This is a reuse of an already analyzed benchmark, not a newly discovered
number-theoretic result \cite{synA,synB}.

\begin{proposition}[Balance before inversion]\label{prop:balance}
If \(d,q,n\in\Z\) satisfy \(dq=n\), every unital ring homomorphism
\(j:\Z\to R\) gives \(j(d)j(q)=j(n)\). When \(R\) is a field and
\(j(d)\ne0\), it follows that \(j(q)=j(n)/j(d)\).
\end{proposition}
\begin{proof}
Apply \(j\) to the balance equation and use preservation of multiplication.
In the field case multiply by the inverse of \(j(d)\), or apply the
corresponding cancellation theorem.
\end{proof}

The type distinction is essential. Nonzero integers may map to zero in positive
characteristic. A regular element can be cancelled but need not be a unit;
\(2\) in \(\Z\) is the elementary example. Sorrel records exact quotient
construction, balance transport, and inversion as different methods.

For \((a,b,p)=(3,2,5)\), the coefficients are \(-53\) and \(-486\).
The executed model checks these values and the four single-premise neighbors
\((3,2,4)\), \((3,3,5)\), \((3,2,3)\), and \((1,2,5)\). They respectively
detect removal of oddness, evenness, the exponent lower bound for \(a_4\),
and the residue condition for \(a_2\). These finite checks are not a replacement
for the universal proof; they are useful tests of the intended benchmark.

The full counterexample package is eventually contradicted in an FLT proof.
That is legitimate mathematics. But testing an arithmetic interface only in
such a context may allow an irrelevant contradiction proof to pass. Sorrel's
strict arithmetic plan uses a restricted helper context with a positive
fixture. It does not outlaw proof by contradiction throughout the language.

\subsection{Sharp regularity is a stronger contract than a bound}
ProveIt's regularity file uses
\[
 F'(x)=2\,\mathrm{up}(2x-1),\qquad 0\le\mathrm{up}\le1,
\]
to establish a global Lipschitz constant two. It then uses
\(F'(1/2)=2\) to show that no smaller constant works \cite{regularity}.
A suitable compact plan is:
\begin{lstlisting}
have abs(F'(x)) <= 2 for all x by differential_law and range_bounds
have Lipschitz(F,2) by mean_value
have F'(1/2) = 2 by differential_law and normalization
have 2 is the least Lipschitz constant of F
  by derivative_lower_bound and the established upper bound
\end{lstlisting}

The third line supplies additional evidence, not decorative explanation.
From \(F\) being \(K\)-Lipschitz, each difference quotient has magnitude
at most \(K\); its limit gives \(|F'(1/2)|\le K\), so \(2\le K\).
Without such a lower-bound argument, the leastness obligation stays open.
Attainment at a point is one sufficient route, not the only possible one.

The source already contains short consequences obtained from the Lipschitz
statement. Sorrel should not replace those with a larger ritual. This example
is a regression test for result strength and a baseline for whether a new
property interface improves on an existing theorem application.

\subsection{A representing algebra needs a dependent package}
The inspected FLT solution assumes a finite index type \(I\), a commutative
ring \(A\), and commutative \(A\)-algebras \(H_i\) that are each finite and
free as \(A\)-modules. It constructs an \(A\)-algebra \(W\), also finite and
free, with equivalences
\[
 e_D:\operatorname{AlgHom}_A(W,D)
       \simeq\prod_{i\in I}\operatorname{AlgHom}_A(H_i,D),
\]
natural in the target algebra \(D\) \cite{tensor}.

Its proof uses induction on a finite type: the empty case uses \(A\); the
successor case tensors the previous representing algebra with the new factor;
the equivalence case transports the indexing family. The successor assembles
several equivalences and then proves naturality componentwise. A Sorrel plan
can make that construction explicit while delegating routine adapter
composition:
\begin{lstlisting}
construct a representing algebra by finite-family induction
  empty: use A with its canonical A-algebra structure
  extension: tensor the previous witness with the new factor
  equivalence: transport the family along the given equivalence
retain the chosen algebra structures and the induced hom equivalences
establish naturality by composition of the construction's naturality laws
\end{lstlisting}

Each step returns data and laws about that data. The finite/free guarantees
must remain attached to the selected algebra and module structures. An
unordered union of evidence from two different constructed witnesses does not
make one package.

The factor hypotheses are necessary for the advertised general guarantee:
for a singleton family, a natural representing object is isomorphic to its
single factor; \(\Q[X]\) is not finite-dimensional over \(\Q\). This is the
premise omission corrected in Moraine's exposition by the syntheses, not a
bug in the original Lean theorem \cite{synA,synB,tensor}.

Finally, the source conclusion is existential. It does not promise an
executable algorithm producing the package from a proof of abstract finiteness.
Sorrel must distinguish choosing a witness inside a proof, declaring a
noncomputable construction under a choice policy, and exporting executable
data. Merely replacing \code{exists} with \code{construct} cannot strengthen
that guarantee.

\section{Leant as a plan-completion service}\label{sec:leant}

\subsection{Preserve the existing acceptance boundary}
Leant's inspected \code{Verified} wrapper explicitly records acceptance by a
supplied verification callback. Its nominal role and opaque constructor protect
association with the candidate; they do not turn that receipt into a kernel
proof or behavioral theorem. The source says that stronger authority must be
retained and replayed separately \cite{leantVerify}.

The interface groups rendered variants, tries them in order, and records the
first accepted representative. Its keyed version deduplicates accepted
renderings without transferring semantic authority to a different occurrence.
Sorrel should retain this careful candidate identity rather than replacing it
with a vague ``AI verified'' status \cite{leantVerify}.

The revised syntheses report a later Leant path that checks a supplied exact
behavioral proposition about a candidate and treats a successful negation
check differently from failure of a positive attempt. Sorrel's integration
must accommodate that stronger baseline while still retaining the proof and
source correspondence needed by its own publication boundary
\cite{synA,synB}.

\subsection{Send a typed residual request, not prose}
A plan-completion request contains the frozen target, the local dependent
context, available premises, permitted providers, and the authorized kind of
hole. It may ask for a proof of a guard, construction of an adapter, or a
function satisfying a behavior specification. A data hole is accompanied by
its specification; the service is not free to alter it.

An adapter is a particularly plausible synthesis target. For example, given
an invariant predicate and a map preserving it, synthesize the map on the
subtype, then prove its behavior from the original law. Another target composes
two contracts using Proposition~\ref{prop:compose}. These exercises genuinely
require term assembly; a successful lookup of one existing theorem should be
reported separately.

A practical implementation can treat difficult dependent propositions as
opaque proof inputs to a structural synthesis fragment, then reconstruct and
check the complete term in Lean. It should not claim that Leant's structural
search automatically decides arbitrary dependent behavior. Unsupported
translations must remain explicit, consistent with Leant's documented
fragment boundaries \cite{leantReadme}.

\subsection{Behavioral pruning before completion}
Suppose a partial program \(s\) denotes a set \(\mathcal C(s)\) of authorized
completions. Rejecting one completion \(c\in\mathcal C(s)\) excludes that
candidate only. Pruning the whole sketch requires a necessary-condition theorem
covering every permitted completion, or a heuristic status that does not claim
logical impossibility.

Residual interfaces offer a disciplined source of such constraints. A consumer
may demand a universally length-preserving operation, an invariant-preserving
map, or a relational pair of outputs. The synthesis engine can propagate these
specifications into holes using checked implication rules. If a propagation
rule only supplies a sufficient construction strategy, failure of that strategy
cannot refute the sketch.

For sorting, the frozen specification should say
\[
 \forall xs,\ \mathsf{Sorted}(f(xs))\land\mathsf{Perm}(f(xs),xs),
\]
with the order relation fixed. Stability is a separate further property if
requested. Length preservation is a useful filter, not a substitute for the
specification. The exact selected function term must occur in the theorem
accepted at the boundary.

\subsection{Positive and negative abstraction contracts differ}
Let an admissible input \(x:A\) have an observation \(o(x):U\), and let a
model \(\widehat f\) commute with the concrete operation through a proved
semantic bridge. A universal model theorem over all \(u:U\), together with a
forward transfer into the desired concrete specification, can prove every
admissible concrete case. Extra unrealizable abstract inputs do not invalidate
such a positive proof.

A negative abstract witness is different. To refute the concrete specification,
there must be an admissible \(x\) realizing the witness and a theorem that a
true concrete specification would imply the tested abstract condition. This
is not implied by positive transfer alone. Surjectivity of the observation is
sufficient but stronger than necessary; realization of the particular witness
suffices.

For list lengths, the possible observations depend on the element type.
At \code{List Empty}, only zero is realized. A constant-empty function and the
identity therefore have the same concrete length behavior, even though their
unrestricted affine models \(0\) and \(n\) differ. With a supplied element,
arbitrary finite lengths are realizable. Two individually realizable
observations also need not be jointly realizable from the same input. Sorrel
retains the entire configuration and its admissibility conditions
\cite{synA,synB}.

The interface consequently distinguishes: a checked source proof; a checked
source refutation; a model-relative counterexample awaiting realization; an
unsupported fragment; and an exhausted search. Only the first two establish a
logical verdict about the original candidate and specification.

\subsection{Replay and retention}
A successful suggestion is accepted by replaying the exact selected source in
the original proof state, accounting for all generated goals. The returned
proof, fixed target, provider identities, environment pin, and axiom policy
are retained. A later pretty-printer rewrite is a new candidate unless it is
rechecked or connected by an adequate typed bridge.

A service can use language models to rank plans, propose intermediate lemmas,
or suggest a construction. Those proposals are untrusted inputs to the same
acceptance boundary. A fluent explanation is not behavioral evidence, and a
model's choice of a more convenient target is not a proof of the author's
request.

\section{Checked computation as an operation with a contract}\label{sec:cas}

\subsection{One request, several permissible execution lanes}
A CAS-style request in Sorrel has an input, an output type, and a fixed
relation connecting the output to the original mathematical object:
\[
 \mathsf{Request}(x):=\Sigma y:O(x),\ \mathsf{Spec}(x,y).
\]
For executable production this is a data package (implemented with an ordinary
structure or subtype when the second field is propositional); a theorem asserting its
existence is a different artifact. Suitable requests include an exact
polynomial identity, a conditional rational-function simplification, an
isolating interval for a specified root, a finite jet of a selected formal
solution, or an enclosure with a stated error. These are separate relations,
not success levels on one scale.

Three lanes can satisfy the request. A verified algorithm has a correctness
theorem and is evaluated through an accepted route. An untrusted searcher
returns a certificate checked by a verified checker. An existing proof-producing
Lean tactic returns a proof directly. The source contract is the same;
checking cost and retained evidence differ. The two syntheses already insist
on this separation \cite{synA,synB}.

\subsection{A shared reification interface, not a universal normalizer}
A reusable reifier can target a ring-expression grammar with literals,
variables, addition, negation, and multiplication. For each original Lean
expression \(e\), it returns a syntax tree \(q\), an ordered environment of
opaque atoms \(v\), and a bridge
\[
 \beta_e:\quad e=\llbracket q\rrbracket_v.
\]
Opaque atoms are exact retained terms, not arbitrary printed names. Different
values may share a symbol only if the actual correspondence is justified.
The grammar and denotation specify the coefficient domain and chosen ring
structure.

A polynomial identity checker may then work with normalized coefficients.
An affine inequality checker needs additional ordered semantics and a theorem
about nonnegative combinations. A length abstraction needs a bridge from list
operations to natural-number expressions, not merely a ring reifier.
Sharing expression syntax does not make these semantic theorems identical.

The original-target chain is explicit:
\[
 e \xlongequal{\beta_e}\llbracket q\rrbracket_v
 \xlongequal{\text{checked certificate}}\llbracket q'\rrbracket_v
 \xlongequal{\beta_{e'}^{-1}}e'.
\]
A correct equality about a different list of atoms proves the wrong theorem;
the reification bridges are therefore part of acceptance, not optional
metadata.

\subsection{Soundness, completeness, and execution are separate theorems}
A typical positive theorem is
\[
 \CheckCert(x,c)=\mathsf{true}\Longrightarrow\mathsf{Spec}(x,\mathsf{out}(c)).
\]
It proves every accepted result without requiring the checker to accept all
true specifications. Decision completeness is a separate converse, restricted
to a specified grammar and admissible domain. Negative certification needs its
own theorem, not the failure of this implication's antecedent.

There is also an execution obligation. A theorem about the checker as a Lean
function does not certify arbitrary bytes reported by a native executable.
The strict publication profile uses a retained proof or a checker computation
replayed through an accepted kernel-level route. Accelerated native evaluation
must be documented under its actual trust and axiom policy for the selected
toolchain. This paper does not claim a performance threshold at which that
tradeoff becomes worthwhile.

\subsection{Precision is an index with algebraic laws}
A jet can be specified by agreement of coefficients below order \(N\), or
congruence modulo \(X^N\). Differentiation reduces a guaranteed order by one:
for \(N\ge1\), agreement modulo \(X^N\) implies derivative agreement modulo
\(X^{N-1}\). Multiplication, composition, and inversion have their own
premises and precision rules. A residual equation for a formal root also
needs a branch or constant-term condition and an applicable uniqueness theorem.

These requirements fit the same proof-plan interface. The consumer states its
needed precision; registered theorems propagate sufficient requirements back
to its inputs. An approximation is never silently presented as equality of
whole functions, and a polynomial root certificate is never silently promoted
to completeness of the solution set. To claim complete solutions, a result
needs both soundness of its listed solutions and coverage of all admissible
solutions.

\section{Implementation architecture and its first useful slice}\label{sec:implementation}

\subsection{Components with explicit trust boundaries}
The proposed Lean package has six components. A source elaborator creates the
frozen header and typed plan nodes. A theorem registry validates signatures
and binder roles. A local property index points to accepted expressions in the
Lean context. A grounding and support service schedules eligible proof routes.
Provider adapters call ordinary tactics, Leant, or certificate generators.
Finally, an evidence compiler and renderer retain the accepted terms and
produce the human-facing explanation.

Only Lean's proof checking determines theorem acceptance. The frontend and
renderer still bear separate responsibilities for source fidelity; their bugs
can make a correct theorem look like a different authored claim. Method-policy
checking adds another obligation and is not folded into a misleading notion
that every trusted function has the same purpose.

The first implementation should generate explicit Lean proof applications and
share intermediate results through local declarations. It need not maintain
a second independently authoritative context or a global database of every
mathematical fact. Provenance and source positions index existing evidence.
The publication artifact includes a reproducible Lean file, not only an editor
screenshot or a process-local receipt.

\subsection{A concrete incremental delivery plan}
\textbf{Slice A: one proof-implicit use site.} Wrap exact quotient transport
with fixed data arguments and explicit divisibility/nonzero premises. Register
only the few arithmetic bridge theorems needed by the nonvacuous fixture.
Accept an existing Lean proof as a way to fill every obligation. The interface
must export the same target under an explicit support-restricted context.

\textbf{Slice B: two sufficient routes.} Add the finite support engine, using
Lean expressions rather than the companion's atom names. Show the direct
and parity-based divisibility routes, remove oddness, and display the
resulting alternative frontier. Require an actual proof for the chosen repair.

\textbf{Slice C: analysis transfer.} Add the Cauchy example's consumers,
including almost-everywhere replacement. Compare with the same helpers and
\code{fun_prop} available in ordinary Lean. Verify that mass two, a point
atom, and a changed measure receive the intended meanings.

\textbf{Slice D: Leant and dependent construction.} Complete a genuine adapter
hole, then inspect the induced-subtype or tensor-family example. Record when
the grounding fragment refuses a data-dependent obligation. That refusal is
a useful boundary, not a reason to approximate dependent types unsafely.

Input grammar beyond the proof-implicit binder is optional until these slices
show a benefit. The full narrative surface can first exist as a read-only
renderer over accepted Lean plans. A renderer-only or library-only outcome is
an acceptable result of the experiment.

\subsection{Lexical structures: an experiment, not a magic preamble eraser}
A lexical structure context records which ring, topology, scalar action, and
measure each expression uses. It helps prevent ambiguity and provides explicit
adapters when an API requires local typeclass instances. It does not imply
that every generated FLT preamble disappears. Some preamble choices control
search or resolve overlapping instances, and a new interface can merely move
that cost elsewhere.

The reviewed tensor solution already relies on carefully scoped structures
and introduces finite/free instances from its inductive witness. Sorrel should
preserve their dependency on that witness. A separate controlled experiment
can compare preamble size, elaboration cost, and repair cost with and without
a lexical context mechanism. This paper does not report that experiment.

\section{The executable companion: what it establishes}\label{sec:companion}

\subsection{A narrow frontend that actually runs}
The archive contains \code{sorrel_model.py}, a dependency-free Python program
implementing the finite calculus of Section~\ref{sec:supports}. Its small
\code{.sorrel} input format declares a lexical context, proposition names,
assumption slots, currently available slots, a finite implication table, and
a goal. This is not the mathematical language illustrated in
Section~\ref{sec:surface}.

The quotient input is:
\begin{lstlisting}
context ExactQuotient
atom EvenB LargeP OddP ResidueA Direct4 Dvd4 Dvd16 Both
assumption EvenB LargeP OddP ResidueA Direct4
available EvenB LargeP ResidueA
rule power16: EvenB & LargeP -> Dvd16
rule power4: EvenB & LargeP & OddP & ResidueA -> Dvd4
rule direct: Direct4 -> Dvd4
rule pair: Dvd4 & Dvd16 -> Both
goal Both
\end{lstlisting}

The output gives two minimal supports, the residual alternatives
\(\{\mathsf{Direct4}\}\) and \(\{\mathsf{OddP}\}\), and the status
\code{open}. The model does not claim to have proved the arithmetic rules.
They are abstract interfaces assumed by this experiment. The separately
executed integer fixtures test selected positive and negative values.

For each retained support, the solver retains a finite proof tree. A separate
checker validates the target, rule identifier, ordered premise list, arity,
context, and exact assumption leaves. It does not consult the solver's labels.
The exporter turns one selected tree into an ordinary Lean implication whose
proposition meanings and used rules remain explicit parameters.

\subsection{Executed results}
The test run used Python 3.13.5 and deterministic seed \(20260906\).
The exact result file is \code{test_results.json}.
\begin{center}\small
\begin{tabularx}{\textwidth}{@{}Y r@{}}\toprule
Executed quantity & Count\\\midrule
Named regression tests & 31\\
Generated finite rule systems & 80\\
Assumption subsets exhaustively checked across those systems & 1,280\\
Subset/goal-atom reachability comparisons & 7,680\\
Retained certificates checked in generated systems & 517\\
Boolean valuations enumerated & 5,120\\
Support soundness checks inside valuations satisfying the rule table & 4,550\\\bottomrule
\end{tabularx}
\end{center}

Each generated system has six proposition atoms, four possible assumption
slots, and ten randomly generated rules of arity zero through three.
For every one of its 16 assumption subsets, an independently implemented
Boolean Horn closure is compared with the support engine's predictions for
all six atoms. All 64 Boolean valuations are also enumerated; valuations
satisfying the rule implications are used to validate the returned supports.
These are finite exhaustive checks of generated cases, not formal proofs of
the Python implementation. The general arguments are the written theorems
in Section~\ref{sec:supports}.

The named tests include alternative supports, residual minimization, closing
and reopening a route, an unseeded cycle, a seeded finite derivation, deletion
of its seed, zero-premise rules, policy filtering, rule-order invariance,
altered targets, sibling contexts, changed rule tables, malformed certificates,
parser errors, exhausted budgets, and the four Frey mutations. A small witness
interface also rejects positive list length when no element witness is
admitted; it is not a Lean proof about \code{Empty}.

\subsection{What was not executed}
No Lean executable was available in this environment. Consequently
\code{SelectedPlan.lean} is labeled \emph{generated but uncompiled}; it is not
counted among checked theorems. The archive does not claim a real Lean
metavariable freeze, typeclass registry, reifier, theorem transport, Leant call,
Cauchy theorem compilation, or benchmark of the repositories.

The miniature parser operates on named propositional atoms, not on Lean
mathematical terms. Its context identifier prevents an elementary scope mixup,
but does not validate equivalence of rebuilt Lean contexts. Its proof checker
checks conditional implication syntax, not the truth of the user-declared
rules. The explicit budget limits rule-combination attempts, not all possible
memory consumption or external provider time. These limits are part of the
artifact's meaning.

The model nevertheless improves on a nonexecuted pseudocode specification:
it fixes input syntax, distinguishes open and supported results, retains
alternative supports, exercises corruption and deletion behavior, and exports
one concrete conditional proof expression. Those are the pieces proposed for
porting to the first Lean slice.

\section{Evaluation: separate the property layer from its helpers}\label{sec:evaluation}

\subsection{Matched arms}
The crucial comparison is not Sorrel against an unassisted author typing raw
Lean from scratch. Use matched arms with shared libraries and the same
mathematical interfaces:
\begin{center}\small
\begin{tabularx}{\textwidth}{@{}P{15mm}Y@{}}\toprule
Arm & Capability\\\midrule
L0 & Ordinary Lean and its applicable existing automation.\\
L1 & L0 plus the newly packaged mathematical helper theorems and adapters.\\
L2 & L1 plus the same Leant, search, and certificate-producing services.\\
L3 & L2 plus Sorrel's proof-implicit interface, property index, and residual plans.\\
L4 & L3 plus the proposed narrative source syntax.\\\bottomrule
\end{tabularx}
\end{center}
A read-only renderer over accepted L2 proofs is an additional useful control.
L3 versus L2 isolates the property/plan layer; L4 versus L3 tests whether new
input syntax is justified. Registration and helper-writing time must be
charged to the arms that need it, rather than treated as free preparatory work.

\subsection{Tasks and intervention logs}
The source files discussed here are development material, not held-out tasks.
They establish acceptance tests and guide implementation. A separate assessor
should select transfer problems after the interface and registry policy are
frozen. Suitable families include local analytic reasoning, invariant-subtype
adapters, and dependent naturality constructions not used while designing the
rules.

Record author interventions, elapsed authoring and repair time, successful
single-lemma lookup versus genuine composition, rule instances generated and
used, support-label sizes, proof-node counts, elaboration/checking time, and
cold versus warm rebuild cost. Retain failures as well as successes. Report
time spent registering a special theorem for one task separately from reused
infrastructure.

An edit study is particularly appropriate for Sorrel. Delete or weaken a
premise, change a measure, rename a binder, change an index, or restrict the
allowed method. Measure whether the fixed target remains unchanged, which
plans survive, which residual obligations appear, and how much author effort
is required to complete a valid repair.

\subsection{Reading comprehension and semantic mutations}
Before revealing an expanded proof, ask readers to reconstruct the compact
claim: objects, domain, quantifiers, measure, branch, exactness, and open
conditions. Include polished invalid neighbors. Then reveal the expansion
and measure whether it corrects misunderstanding. A compact source that
reliably makes readers believe a stronger theorem is not successful compression.

The initial acceptance suite should include at least the following paired
situations. Each negative case needs a nearby admissible positive fixture;
some negative outcomes are rejected promotions, not false mathematical
statements.
\begin{center}\small
\begin{longtable}{@{}P{35mm}P{59mm}P{49mm}@{}}\toprule
Boundary & Positive fixture & Required response to mutation\\\midrule\endhead
Conditional evidence & All leaves of a supported plan are supplied. & Removing one leaf reopens the plan; no silent extra premise.\\
Locality & Equality on a neighborhood of the evaluation point. & Equality only at the point does not justify derivative transport.\\
Almost everywhere & \(\mu\)-a.e. equality used by a \(\mu\)-integral consumer. & Changing \(\mu\) or requesting a point value requires new evidence.\\
Exact division & The satisfiable fixture \((3,2,5)\). & The four removed-premise fixtures fail the corresponding divisibility claim.\\
Result strength & Lipschitz upper bound plus a lower-bound argument. & An upper bound alone does not establish leastness.\\
Abstract refutation & A concrete input realizes the failing observation. & Positive length with no element witness remains unrealized.\\
Structure choice & Finite/free factor algebras and retained selected structures. & Finite indexing alone does not supply finite/free factors.\\
Method fidelity & A permitted rule tree with restricted helper premises. & An unrelated contradiction proof does not satisfy that method policy.\\
Source integrity & Every displayed assertion has accepted evidence. & An unused false assertion prevents complete publication.\\
Rebuild integrity & Evidence is rechecked in the destination context. & Old local identifiers or matching printed names alone are insufficient.\\\bottomrule
\end{longtable}
\end{center}

\subsection{Stopping rules}
Set quantitative thresholds after a pilot estimates variability, and before
the main comparison. This paper invents no productivity percentage. The
separate-language experiment should stop or shrink when L3's gains vanish
against L2, its explanations increase semantic errors, or its registration
and repair costs exceed the work it saves.

A useful outcome may be a library of mathematical contracts, an evidence-aware
editor, a Leant adapter, or a read-only renderer. Retaining those pieces while
dropping the narrative language would follow the evidence rather than defeat
the proposal.

\section{Answers to the two reports' fourth-iteration questions}\label{sec:crosswalk}
The two question lists overlap but emphasize different implementation risks.
The following crosswalk states Sorrel's decision and the evidence still needed,
rather than treating a design answer as an experiment already completed
\cite{synA,synB}.
\begin{center}\small
\begin{longtable}{@{}P{34mm}P{70mm}P{39mm}@{}}\toprule
Question cluster & Sorrel decision & Current status\\\midrule\endhead
First use site and package & Proof-implicit exact-quotient transport; two alternative divisibility routes; restricted arithmetic premises. & Finite symbolic slice executed; Lean integration proposed.\\
Finite rule generation & Explicit finite term bank, schema whitelist, bounded constructors and instantiations; exactness only relative to generated rules. & Algorithm and bounds specified; Lean grounding not built.\\
Data inference and structures & Freeze meaning-bearing data; delegate witnesses only through a fixed specification; retain dependent structure packages. & Design; source examples inspected.\\
Required proof methods & Filter rules before minimization; check typed plan nodes and support-restricted helper contexts separately from kernel truth. & Rule-filter test executed; full policy compiler proposed.\\
Behavioral bridge and completeness & Keep positive soundness, exact decision, and negative realization as separate provider contracts. & Interface analysis; no new Lean behavioral checker.\\
Negative Leant outcomes & Refutation needs source proof or realized counterexample with transfer; one failed completion cannot refute a sketch. & Source review and design.\\
Nonalgebraic transfer and a.e. consumers & Cauchy--CDF example, measure-anchored relation, explicit integral-congruence consumer. & Mathematical reconstruction; not a held-out task.\\
Rewriting, indices, and lifetime & Keep original anchors; insert checked bridges; rebuild per declaration before attempting cache transport. & Elementary context rejection tested; dependent transport proposed.\\
Reifier and checker reuse & Shared ring syntax with per-expression denotation proofs; distinct ordered and behavioral semantic layers. & Specification only.\\
Actual cost and rule usage & Instrument labels, generation, lookup/composition, proof growth, and cold/warm elaboration in matched arms. & Finite operation counts only; no Lean benchmark.\\
Lexical preambles & Measure instance-context changes rather than assume a preamble disappears. & No preamble-compression result claimed.\\
Reader recovery and number of implementations & Blind reading and edit studies; independently implement/check the same small interface before expanding syntax. & Evaluation protocol; no participants tested.\\
Evidence migration and stop conditions & Recheck exact target/dependencies at destination; keep library/services if the language adds no benefit. & Acceptance policy and proposed stopping rule.\\\bottomrule
\end{longtable}
\end{center}

\section{Risks, priorities, and conclusion}\label{sec:conclusion}

\subsection{The largest risks}
The first risk is search explosion. Minimal support sets make explanations
more informative but can become exponentially numerous. The initial engine
therefore needs small profiles, explicit completeness status, and demand-driven
use. Top-ranked partial explanations can be useful, but must not be advertised
as the complete set of repairs.

The second risk is an expensive second elaborator. Duplicating instance search,
function-property automation, and rewriting can erase the author's savings.
Sorrel should orchestrate existing proof-producing tools and retain only
additional information that is demonstrably useful for source meaning,
residual interfaces, or repair.

The third risk is false precision in explanations. A minimal support in a
fragment can still be redundant in mathematics, and a strict method grammar
can reject perfectly reasonable alternative proofs. The interface must state
its scope and permit an explicit author choice to change the proof route.

The fourth risk is representational drift: changed measures, coefficient
structures, universes, witnesses, and dependent indices. An uncompiled
finite prototype does not solve it. Expression-level acceptance checks and
rebuild tests are therefore higher priority than more surface vocabulary.

\subsection{What to implement next}
The next engineering artifact should be the exact-quotient use site in Lean,
with the source target frozen, one property-implicit binder, the two support
routes, explicit residual display, and a retained kernel-checked proof. The
independent evaluator should receive the paired mutations before the frontend
is expanded. The Cauchy example should follow as an analysis transfer test
within the development suite, not as a claimed unseen benchmark.

Only after those experiments should Sorrel's narrative input be enlarged or
a native refinement former be attempted. The mathematical property types are
important; their successful use in an actual proof state is the decisive unit.

\subsection{Conclusion}
Sorrel proposes a type system for \emph{arguments in progress}. An object's
proved properties determine which operations can consume it. A function's
behavior is a theorem about its actual evaluations. A mathematical plan
states its target and exports a typed interface of the evidence still needed.
Alternative sufficient interfaces can be retained and inspected without
pretending to compute globally weakest assumptions.

This changes what authors must write while preserving the role of Lean's
kernel. The author selects the mathematical route and meaningful construction
choices; Lean checks the completed evidence; Leant and certified computation
fill permitted gaps; the interface explains remaining conditions without
changing the proposition.

The contribution is deliberately narrower than a finished language. The
finite calculus has written proofs and an executed reference model. The
corpus reconstructions identify what a real frontend must preserve. The
measurement protocol can reject the separate-language hypothesis while keeping
useful libraries and tools. That is a stronger fourth-round position than
another promise that better syntax alone will make formal mathematics concise.

\clearpage\appendix
\section{Reproducing the companion and the article}\label{app:repro}
The archive is self-contained for rebuilding the PDF with a conventional
pdfLaTeX installation. The article embeds its bibliography and requires no
external figures or font files. The companion requires Python 3.10 or newer
and no third-party packages.
\begin{lstlisting}
python sorrel_model.py --test --json test_results.json
python sorrel_model.py examples/quotient.sorrel \
  --json examples/quotient_result.json \
  --export-lean SelectedPlan.lean
python sorrel_model.py examples/cauchy.sorrel \
  --json examples/cauchy_result.json
pdflatex -interaction=nonstopmode -halt-on-error sorrel.tex
pdflatex -interaction=nonstopmode -halt-on-error sorrel.tex
pdflatex -interaction=nonstopmode -halt-on-error sorrel.tex
\end{lstlisting}

A budgeted model run is requested with \code{--budget N}. An incomplete support
inventory can still contain a valid conditional derivation; the output exposes
both statuses. A malformed input exits with an error rather than silently
inventing a proposition or rule.

The generated Lean specimen uses \code{Init} and states a propositional
implication with its used rule interfaces as hypotheses. Running
\code{lean SelectedPlan.lean} is an intended additional check on a machine
with Lean installed. It was not performed here, and compiling it would still
not prove the mathematical meaning of the model's rule names.

\section{A compact interface schema}\label{app:schema}
The following is a design schema, not a claimed implementation:
\begin{lstlisting}
FrozenRequest
  target        : elaborated Lean expression
  context       : ordered local declaration telescope
  universes     : fixed level assignments
  protectedData : fixed structures, domains, witnesses, and relations
  allowedHoles  : proof goals or explicitly specified constructions

PlanNode
  originalRequest : FrozenRequest
  operation       : checked theorem interface or explicit constructor
  dataArguments   : typed expressions
  premises        : ordered premise nodes
  residual        : well-formed dependent telescope
  result          : expression checked under the residual telescope
  methodPolicy    : optional accepted-plan grammar and dependency policy

AcceptedArtifact
  statement      : exact checked target
  evidence       : retained Lean term or replayable checked certificate
  sourceMap      : every displayed assertion to its evidence node
  dependencies   : environment and actual trust/axiom policy
  planRecord     : required method evidence, separately validated
\end{lstlisting}

Publication checks all display nodes and confirms that each frozen request's
residual telescope is filled. A user may instead publish an explicitly
conditional theorem; its statement then contains the conditions. A conditional
artifact and an unconditional theorem must not share a status that suppresses
this difference.

\section{Provenance and bibliographic notes}\label{app:sources}
Repository citations below link to the pinned files, not to mutable default
branches. The record \code{sources.json} includes complete commit and blob
identities for the main inspected files. Line ranges stated in that record
are source-file ranges, not PDF page numbers. No upstream source archive is
redistributed. The inspection date for live official documentation is
6 September 2026.

The reports' attributions to earlier proposals are retained as attributions:
Trellis's finite qualifier closure motivates making the grounding boundary
explicit; Fiber's and Obsidian's nonvacuity warnings motivate the arithmetic
fixture; Moraine's behavioral model and Karst's empty-type example motivate
separate realization contracts. Sorrel's finite support sets also have a clear
prior-art relationship to the ATMS label of minimal environments. Unlike a
full default-reasoning ATMS, this model neither assumes consistency checking
nor removes mathematically contradictory contexts automatically. It computes
positive conditional derivability from a fixed implication table.

\begin{thebibliography}{99}\small\raggedright
\bibitem{synA} Codex.
\emph{From Properties to Proof Obligations: Round 3 Synthesis}.
Brainstorm, revised round-three synthesis, 6 September 2026.
Commit \code{bf3333905995060870f8585e297ffc16f3fe7e86}.
\url{https://github.com/VladimirReshetnikov/Brainstorm/blob/bf3333905995060870f8585e297ffc16f3fe7e86/docs/round-3/synthesis/unified-report.tex}.

\bibitem{synB} Claude (Anthropic).
\emph{Nine Refinements: Property-Aware Typing in a Mathematician's Proof Language over Lean}.
Brainstorm, revised round-three review, 6 September 2026.
Same Brainstorm commit as above.
\url{https://github.com/VladimirReshetnikov/Brainstorm/blob/bf3333905995060870f8585e297ffc16f3fe7e86/docs/round-3/unified_report/unified_report.tex}.

\bibitem{cauchy} ProveIt contributors.
\emph{Cauchy transforms through compactly supported CDFs}.
\code{Analysis/FabiusFunction/Lean/FabiusFunction/CauchyCDF.lean},
commit \code{24ce8bd743eaab64a91ce90725ea00f498d319d2}.
\url{https://github.com/VladimirReshetnikov/ProveIt/blob/24ce8bd743eaab64a91ce90725ea00f498d319d2/Analysis/FabiusFunction/Lean/FabiusFunction/CauchyCDF.lean}.

\bibitem{regularity} ProveIt contributors.
\emph{Sharp Lipschitz regularity of the Fabius and Rvachev functions}.
\code{Analysis/FabiusFunction/Lean/FabiusFunction/Regularity.lean},
same ProveIt commit.
\url{https://github.com/VladimirReshetnikov/ProveIt/blob/24ce8bd743eaab64a91ce90725ea00f498d319d2/Analysis/FabiusFunction/Lean/FabiusFunction/Regularity.lean}.

\bibitem{frey} Fermat repository contributors.
\code{Definitions/Def_FLTPrelim_FreyPackage.lean}.
Source credits Kevin Buzzard, Ruben Van de Velde, and Pietro Monticone;
ported from the Imperial College London FLT formalization.
Commit \code{aa2d8b34692b16c70f699536de0d8e75b9a3e9ef}.
\url{https://github.com/anthropics/fermats-last-theorem/blob/aa2d8b34692b16c70f699536de0d8e75b9a3e9ef/Definitions/Def_FLTPrelim_FreyPackage.lean}.

\bibitem{tensor} Fermat repository contributors.
\code{P2M/Sol/S_Algebra_exists_algHom_equiv_pi.lean}.
Same Fermat repository commit; finite/free factor hypotheses in the
\code{solution} theorem, followed by the naturality conclusion and proof.
\url{https://github.com/anthropics/fermats-last-theorem/blob/aa2d8b34692b16c70f699536de0d8e75b9a3e9ef/P2M/Sol/S_Algebra_exists_algHom_equiv_pi.lean}.

\bibitem{leantVerify} Vladimir Reshetnikov and Leant contributors.
\code{src/Leant/Synth/Verification.hs}.
Commit \code{3a40904be8a410d832d9ab6900b3d3e7b425eccb}.
\url{https://github.com/VladimirReshetnikov/Leant/blob/3a40904be8a410d832d9ab6900b3d3e7b425eccb/src/Leant/Synth/Verification.hs}.

\bibitem{leantReadme} Vladimir Reshetnikov and Leant contributors.
\emph{Leant: interactive Lean 4 helper}, repository README.
Public description inspected 6 September 2026; implementation claims in this
article are bounded by the specific source review noted above.
\url{https://github.com/VladimirReshetnikov/Leant}.

\bibitem{funprop} Mathlib contributors.
\emph{Mathlib.Tactic.FunProp}, official module documentation.
\url{https://leanprover-community.github.io/mathlib4_docs/Mathlib/Tactic/FunProp.html}.

\bibitem{leanSubtype} Lean contributors.
\emph{The Lean Language Reference: Subtypes}.
\url{https://lean-lang.org/doc/reference/latest/Basic-Types/Subtypes/}.

\bibitem{leanAxioms} Lean contributors.
\emph{The Lean Language Reference: Axioms}; and
\emph{Theorem Proving in Lean 4: Propositions and Proofs}.
\url{https://lean-lang.org/doc/reference/latest/Axioms/}.
\url{https://lean-lang.org/theorem_proving_in_lean4/Propositions-and-Proofs/}.

\bibitem{lean433} Lean contributors.
\emph{Lean 4.33.0 release notes}, 10 August 2026.
In particular, automation, incremental elaboration, and the renaming of
experimental \code{mvcgen'} to \code{vcgen}, leaving \code{mvcgen} unchanged.
\url{https://lean-lang.org/doc/reference/latest/releases/v4.33.0/}.

\bibitem{dekleer} Johan de Kleer.
\emph{Problem Solving with the ATMS}.
Artificial Intelligence 28 (1986), 197--224.
The introductory description of minimal-environment labels and implication
justifications is especially relevant.
\url{https://dekleer.org/Publications/Problem%20Solving%20with%20the%20ATMS.pdf}.

\bibitem{provenance} Todd J. Green, Grigoris Karvounarakis, and Val Tannen.
\emph{Provenance Semirings}.
Proceedings of PODS 2007, 31--40.
\url{https://doi.org/10.1145/1265530.1265535}.

\bibitem{lovas} William Lovas and Frank Pfenning.
\emph{Refinement Types for Logical Frameworks and Their Interpretation as
Proof Irrelevance}.
Logical Methods in Computer Science 6(4), 2010; arXiv:1009.1861.
\url{https://arxiv.org/abs/1009.1861}.

\bibitem{lean4lean} Mario Carneiro.
\emph{Lean4Lean: Verifying a Typechecker for Lean, in Lean}.
arXiv:2403.14064, 2024.
\url{https://arxiv.org/abs/2403.14064}.
\end{thebibliography}
\end{document}
